\documentclass[11pt]{article}
% Shared design for the two minimal papers.
% It follows the layout of the guide "The Riemann Hypothesis and its generalizations"
% (KNT with Fable, 9 November 2026): a designed title block, a short story list at the start,
% coloured headings, boxed statements, idea boxes, and annotated figures.
%
% A paper sets \papertitle, \papersubtitle and \paperrunning before \begin{document}
% and calls \maketitleblock after it.

\usepackage[letterpaper,left=1.12in,right=1.12in,top=1.05in,bottom=1.1in]{geometry}
\usepackage{amsmath,mathtools,amssymb,amsthm}
\usepackage{newpxtext}
\usepackage[varbb]{newpxmath}
\usepackage{fontspec}
\setsansfont{texgyreheros}[Extension=.otf, UprightFont=*-regular, BoldFont=*-bold,
  ItalicFont=*-italic, BoldItalicFont=*-bolditalic, Scale=0.88]
\usepackage{microtype}
\usepackage{parskip}
\usepackage{xcolor}
\usepackage{booktabs,array}
\usepackage{enumitem}
\usepackage{multicol}
\usepackage[titles]{tocloft}
\usepackage{tikz}
\usetikzlibrary{decorations.markings,arrows.meta,positioning,calc}
\usepackage[most]{tcolorbox}
\usepackage{fancyhdr}
\usepackage{titlesec}
\usepackage[font=small,labelfont={bf,sf},labelsep=period,skip=6pt]{caption}
\usepackage{url}

\definecolor{accent}{HTML}{1B4F72}
\definecolor{warm}{HTML}{C0562A}
\definecolor{slate}{HTML}{56636E}
\definecolor{ink}{HTML}{1E1E1E}

\usepackage[colorlinks=true,linkcolor=accent,citecolor=accent,urlcolor=accent]{hyperref}

% Lean identifiers: typewriter, breakable at dots and underscores.
\DeclareUrlCommand\lean{\urlstyle{tt}}

% ---------- layout ----------
\linespread{1.06}
\setlength{\parskip}{0.55em plus 0.1em minus 0.05em}
\setlist{itemsep=0.25em, topsep=0.2em, parsep=0.1em}
\titleformat{\section}{\sffamily\bfseries\Large\color{accent}}{\thesection}{0.7em}{}
\titleformat{\subsection}{\sffamily\bfseries\normalsize\color{ink}}{\thesubsection}{0.6em}{}
\titleformat{\paragraph}[runin]{\sffamily\bfseries\small\color{accent}}{}{0pt}{}[\enspace]
\titlespacing*{\section}{0pt}{22pt}{5pt}
\titlespacing*{\subsection}{0pt}{13pt}{1pt}
\titlespacing*{\paragraph}{0pt}{9pt}{0pt}
\pagestyle{fancy}
\fancyhf{}
\renewcommand{\headrulewidth}{0pt}
\fancyhead[L]{\sffamily\footnotesize\color{slate}\paperrunning}
\fancyhead[R]{\sffamily\footnotesize\color{slate}\thepage}
\setlength{\headheight}{14pt}
\setcounter{tocdepth}{1}
\setlength{\cftbeforesecskip}{2pt}
\renewcommand{\cftsecfont}{\small}
\renewcommand{\cftsecpagefont}{\small\color{slate}}
\renewcommand{\cftsecleader}{\cftdotfill{\cftdotsep}}
\setlength{\cftsecnumwidth}{1.7em}
\setlength{\emergencystretch}{2em}
\numberwithin{equation}{section}

% ---------- title block ----------
\newcommand{\papertitle}{}
\newcommand{\papersubtitle}{}
\newcommand{\paperrunning}{}
\newcommand{\paperauthors}{KNT\enspace\textperiodcentered\enspace Astra\enspace\textperiodcentered\enspace Fable~5.2}
\newcommand{\paperdate}{Draft of 11 November 2026}
\newcommand{\maketitleblock}{%
  \thispagestyle{plain}%
  {\sffamily\bfseries\color{accent}\fontsize{25}{29}\selectfont\raggedright\papertitle\par}%
  \vspace{5pt}%
  {\sffamily\Large\color{slate}\raggedright\papersubtitle\par}%
  \vspace{8pt}%
  {\color{accent}\hrule height 1.2pt}%
  \vspace{7pt}%
  {\sffamily\small\color{ink}\paperauthors\hfill\color{slate}\paperdate\par}%
  \vspace{10pt}}

% Small headings inside the front matter.
\newcommand{\fronthead}[1]{\par\vspace{6pt}{\sffamily\bfseries\normalsize\color{ink}#1}\par\vspace{-2pt}}
\makeatletter
\newcommand{\contentsblock}{%
  \par\noindent\begin{minipage}{\linewidth}%
  \fronthead{Contents}%
  \setlength{\parskip}{0pt}\setlength{\columnsep}{2.2em}\raggedright
  \begin{multicols}{2}
  \@starttoc{toc}
  \end{multicols}
  \end{minipage}\par}
\makeatother

% ---------- statements ----------
% Theorems, lemmas, propositions and corollaries share one counter.
% Each statement is in a tinted box with an accent rule at the left.
\newtheoremstyle{boxed}{0pt}{0pt}{\normalfont}{0pt}{\sffamily\bfseries\small\color{accent}}{.}{0.6em}{}
\theoremstyle{boxed}
\newtheorem{theorem}{Theorem}[section]
\newtheorem{lemma}[theorem]{Lemma}
\newtheorem{proposition}[theorem]{Proposition}
\newtheorem{corollary}[theorem]{Corollary}

\tcbset{guidebox/.style={enhanced, parbox=false, boxrule=0pt, frame hidden,
  arc=1.5pt, left=11pt, right=10pt, top=7pt, bottom=7pt, before skip=11pt, after skip=9pt}}
% Not breakable: a label at the start of a box that breaks at the top of a page can be lost.
\tcbset{statementbox/.style={guidebox, colback=accent!6,
  borderline west={3pt}{0pt}{accent}}}
\tcolorboxenvironment{theorem}{statementbox}
\tcolorboxenvironment{lemma}{statementbox}
\tcolorboxenvironment{proposition}{statementbox}
\tcolorboxenvironment{corollary}{statementbox}

% The main result of a paper: a framed box with a title bar.
% \begin{maintheorem}{name} uses the counter of the theorems, and \label works after it.
\newtcolorbox{mainresult}[1]{enhanced, parbox=false, colback=white, colframe=accent, boxrule=0.9pt,
  arc=1.5pt, left=11pt, right=11pt, top=8pt, bottom=8pt, title={#1},
  fonttitle=\sffamily\bfseries\small, colbacktitle=accent, coltitle=white,
  before skip=12pt, after skip=12pt}
\newenvironment{maintheorem}[1]%
  {\refstepcounter{theorem}\begin{mainresult}{Theorem~\thetheorem\ (#1)}}%
  {\end{mainresult}}

% ---------- proofs ----------
% A proof has a thin rule at the left, so that the eye finds the statements.
\renewcommand{\proofname}{\upshape\sffamily\bfseries\small\color{slate}Proof}
\renewcommand{\qedsymbol}{\textcolor{slate}{\rule{0.55em}{0.55em}}}
\tcolorboxenvironment{proof}{enhanced, breakable, parbox=false, boxrule=0pt, frame hidden,
  colback=white, borderline west={1.2pt}{0pt}{slate!35}, left=11pt, right=0pt, top=1pt, bottom=2pt,
  before skip=6pt, after skip=11pt}
% A step or a case inside a proof.
\newcommand{\step}[1]{{\sffamily\bfseries\small\color{slate}#1}}

% ---------- idea, notation and fine-print boxes ----------
\newtcolorbox{idea}[1]{guidebox, breakable, colback=warm!6, borderline west={3pt}{0pt}{warm!80},
  before upper={{\sffamily\bfseries\footnotesize\color{warm!75!black}IDEA\enspace\textbar\enspace}{\sffamily\bfseries\small #1}\par}}
\newtcolorbox{notation}[1]{guidebox, breakable, colback=black!4, borderline west={3pt}{0pt}{slate!60},
  fontupper=\small,
  before upper={{\sffamily\bfseries\footnotesize\color{slate}NOTATION\enspace\textbar\enspace}{\sffamily\bfseries\small #1}\par}}
\newtcolorbox{fineprint}{guidebox, breakable, colback=white, borderline west={1.5pt}{0pt}{black!25},
  fontupper=\footnotesize\color{black!78}, left=9pt, top=1pt, bottom=1pt,
  before upper={{\sffamily\bfseries Fine print.}\enspace}}

% ---------- figures ----------
\tikzset{
  vtx/.style={circle, fill=ink, inner sep=1.35pt},
  port/.style={circle, draw=accent, fill=white, line width=0.8pt, inner sep=1.5pt},
  edge/.style={line width=0.55pt, draw=ink},
  hot/.style={line width=1.1pt, draw=warm},
  cool/.style={line width=1.1pt, draw=accent},
  region/.style={draw=accent, fill=accent!7, line width=0.6pt, rounded corners=6pt},
  note/.style={font=\sffamily\footnotesize, text=slate},
}

% ---------- macros of the mathematics ----------
\newcommand{\Z}{\mathcal{Z}}
\newcommand{\PP}{\mathbb{P}}
\newcommand{\EE}{\mathbb{E}}
\newcommand{\comp}{\operatorname{comp}}
\newcommand{\tr}{\operatorname{tr}}
\newcommand{\lgk}[1]{\log^{(#1)}}
\newcommand{\lgs}{\log^{*}}
\newcommand{\sm}{\setminus}
\newcommand{\half}{\tfrac12}
\newcommand{\cO}{\mathcal{O}}


\renewcommand{\papertitle}{Pieces at girth five:\\ the case of bounded girth is not necessary}
\renewcommand{\papersubtitle}{A second research note on ``The maximum number of cycles in a graph of large rank''}
\renewcommand{\paperrunning}{Pieces at girth five}
\renewcommand{\paperauthors}{Fable~5.2, for KNT}
\renewcommand{\paperdate}{Research note of 12 November 2026}
\hypersetup{pdftitle={Pieces at girth five: the case of bounded girth is not necessary},pdfauthor={Fable 5.2, for KNT}}

\newcommand{\M}[1]{[M,~#1]}
\newcommand{\N}[1]{[N1,~#1]}

\begin{document}
\maketitleblock

This note answers the second research task of 12 November 2026: can the case of bounded girth
(Section~7 of the paper~[M]) be removed by a stronger lemma on pieces? The answer is yes. The
numbers of the statements of~[M] are those of the draft of 42~pages, and [N1] is the first note,
\texttt{one-scale.pdf}. The manuscript is not changed.

\fronthead{The result in eight sentences}
\begin{enumerate}[leftmargin=1.6em]
\item \emph{The audit.} Theorem~4.4 of~[N1] is correct. Section~\ref{sec:audit} writes its
induction step again, with the two regions and all conditions.
\item \emph{The new theorem.} The lemma on pieces holds at girth five: a connected subcubic
multigraph of girth at least five with at most $k$ boundary vertices and at least $C_1k^{7}$
vertices has a forest proportion of less than $\frac35$ for every even set of boundary vertices
(Theorem~\ref{thm:poly}).
\item It needs no hypothesis on the cuts of an ambient graph. A cubic graph without a cyclic cut
of size at most four has girth at least five, and this is all that the proof uses.
\item \emph{The reason.} A piece with a larger forest proportion is thin. A long thin piece
contains a part in which \emph{every} join is a forest. In such a part, a tree that connects the
terminals passes one time through a region far from the terminals. If the number of terminals on
each side is odd, this region has a full pair of ports. By \M{Lemma~5.8} it is a capped ladder,
and a ladder has cycles of length four.
\item \emph{The consequence.} \M{Theorem~8.1} holds for girth five. Then Case~4 in the proof of
the main theorem is empty. Section~7 (Theorem~7.1), Lemma~5.5 and Corollary~5.6 are not
necessary, and they are used nowhere else (Section~\ref{sec:last}).
\item \emph{Measured.} Two copies of the paper without Section~7 build. The lean form has 1617
typeset lines and 42 pages (now 1624 and 42). The clean form has 1643 lines and 43 pages. The
main theorem has four cases in the place of five.
\item \emph{Tests.} In all families that I computed, the largest forest proportion decreases
with the size, and the ladder is the only family with a rigid set of unbounded size
(Section~\ref{sec:tests}).
\item \emph{The price.} The constant $C_1$ is very large, the proof of the lemma on pieces is
longer, and the text differs more from the formal proof, which has the case of bounded girth.
\end{enumerate}

\begin{notation}{the status of this note}
\begin{itemize}[leftmargin=1.2em]
\item \emph{Established on paper.} Sections~\ref{sec:audit} to~\ref{sec:last} contain complete
proofs. They use \M{Theorem~3.1} through \M{Lemma~4.1}, and \M{Lemmas~5.2 and~5.8}. I checked
each step against these statements.
\item \emph{Tested by computer.} Section~\ref{sec:tests} has complete enumerations on small
graphs. These tests do not replace a proof, and the constants of the theorem are far outside
their range.
\item \emph{Not done.} There is no formal proof. No other person or model has read this note.
\item \emph{The manuscript.} It is unchanged. The script \texttt{apply\_girth\_five.py} applies
either form to a copy.
\end{itemize}
\end{notation}

\contentsblock

\section{The setting}\label{sec:setting}

A graph may have loops and parallel edges; a multigraph of girth at least five has none. Let $N$
be a connected subcubic multigraph. Its vertices of degree less than three are its
\emph{boundary vertices}, and its \emph{deficit} $d(N)$ is the sum of $3-\deg v$ over all
vertices $v$. A count of edge ends gives
\begin{equation}\label{eq:deficit}
  |V(N)|=2\beta(N)-2+d(N).
\end{equation}
Let $T$ be a set of boundary vertices of even size. A \emph{$T$-join} is a set of edges in which
exactly the vertices of $T$ have odd degree. There are $2^{\beta(N)}$ of them. The \emph{forest
proportion} of $T$ is the proportion of the $T$-joins that are forests. We call $T$ \emph{rigid}
if every $T$-join is a forest. For two vertices $x\ne y$ of degree at most two, an
$\{x,y\}$-join is a path from $x$ to $y$ and some cycles, pairwise disjoint; so the set
$\{x,y\}$ is rigid if and only if every $\{x,y\}$-join is a path. This is a \emph{full pair}
in~[M].

\begin{fineprint}
\textbf{Statements that the note uses.}

\smallskip
\textbf{(a) Thin pieces} (\M{Lemma~4.1} and the step ``A marker'' in the proof of
\M{Lemma~8.2}). There is an absolute constant $B>2^{100}$ with the following property. Let $N$
have at most $k$ boundary vertices, let $o$ be one of them, and suppose that some even set of
boundary vertices has forest proportion at least $\frac35$. Then fewer than $W_k=B(k+1)^{2}$
vertices have any given distance from $o$. This step uses no hypothesis on the girth.

\smallskip
\textbf{(b) Poles} (\M{Lemmas~5.2 and~5.8}). A \emph{pole} is a connected multigraph without
loops in which every vertex has degree two or three; the vertices of degree two are its
\emph{ports}. Let $x,y$ be a full pair of ports of a pole $M_1$ of rank $b\ge1$. Replace every
other port, with its two edges, by one edge. The result $M_1^{\circ}$ consists of bridges and of
bridgeless multigraphs $N_1,\dots,N_q$ of ranks $b_i\ge1$ with $b=\sum_ib_i$. Each $N_i$ has
two ports, and they form a full pair. Each $N_i$ is a pair of parallel edges, or a \emph{capped
ladder}: a chain of tiles, where a tile is a copy of $K_{2,2}$ (a \emph{square}, of weight~$2$) or
one edge (a \emph{single rung}, of weight~$1$), two disjoint edges join the right pair of each
tile to the left pair of the next tile, and each of the two ports is joined to the two vertices
of an end pair. A capped ladder of total weight $w$ has rank $w+1$.

\smallskip
\textbf{(c) Separators} (\M{Lemma~8.4}, \N{Lemma~1.2}). Pairwise disjoint connected sets that
separate two vertices $u$ and $p$ of a connected graph are linearly ordered,
$C_1<C_2<\dots<C_n$. For an integer $w\ge1$ there are pairwise disjoint sets $Y_l$, one for each
$l\ge0$ with $(l+1)w\le n-1$, such that $Y_l$ contains $w-1$ of the sets $C_i$, every edge that
leaves $Y_l$ ends in $C_{lw+1}$ or in $C_{(l+1)w+1}$, and $Y_l$ lies in the component of $u$
without $C_{(l+1)w+1}$.

\smallskip
\textbf{(d) Free regions} (\N{Lemma~4.1}). A region $A$ of $N$ is \emph{$T$-free} if it contains
no vertex of $T$ and every component of $N-A$ contains an even number of vertices of $T$. If $N$
has $n$ pairwise disjoint $T$-free cyclic regions with at most $d$ boundary edges each, then the
forest proportion of $T$ is at most $\half+(d+1)2^{d-1}/n$.

\smallskip
\textbf{(e) Product} (\N{Lemma~4.2}, with numbers $f_i$ in the place of $\frac35$; the proof is
the same). Let $U_1,\dots,U_m\subseteq V(N)\sm T$ be pairwise disjoint sets of vertices of
degree three such that each $N[U_i]$ is connected, and let $f_i$ be a number such that every
even set of boundary vertices of $N[U_i]$ has forest proportion at most $f_i$. Then the forest
proportion of $T$ is at most $f_1\cdots f_m$.

\smallskip
\textbf{(f) Few short cycles} (\N{Lemma~4.3}). For $k\ge1$ let $s_k$ be the least integer with
$2^{s_k}>(2s_k+1)W_k$. Let $N$ have at least two and at most $k$ boundary vertices, put $s=s_k$,
and let $t$ be the size of a maximal family of pairwise disjoint cycles of length at most $2s$.
If $N$ has more than $3\cdot2^{s}(k+2s\,t)$ vertices, then every even set of boundary vertices
has forest proportion less than $\frac35$. This lemma has no hypothesis on the girth.
\end{fineprint}

\section{The audit of Theorem 4.4 of [N1]}\label{sec:audit}

\begin{theorem}[{\N{Theorem~4.4}}]\label{thm:old}
Fix $k_0$ such that $4s_k<k$ for all $k\ge k_0$, and put $D_*=2s_{k_0}+1$. There is an absolute
constant $C_1$ with the following property. Let $k\ge1$, and let $N$ be a connected subcubic
multigraph of girth at least $D_*$ with at most $k$ boundary vertices and at least $C_1k^{7}$
vertices. Then every even set of boundary vertices has forest proportion less than $\frac35$.
\end{theorem}

\paragraph{The result of the audit.} The proof is correct. One passage was not precise: the proof
took ``one of the others'' and then used ``two sets $Y'$''. The step is written in full below,
and~[N1] has the new wording. I found no other defect. The proof uses the girth in two places
only. These are the base, and the number of vertices of a short cycle in~(6) below.

\paragraph{The numbers.} For $k<k_0$ put $F(k)=3\cdot2^{s_k}k+1$. For $k\ge k_0$ put $s=s_k$,
$k'=4s<k$ and
\begin{gather*}
  w=\lceil k'F(k')/D_*\rceil+1,\qquad n_1=5(2s+1)\,2^{2s},\qquad n_2=k^{2}(k+2)\,w,\\
  F(k)=3\cdot2^{s}\bigl(k+2s\,(k+n_1+n_2+2)\bigr)+1 .
\end{gather*}
The claim $P(k)$ is: a multigraph $N$ as in the theorem with at least $F(k)$ vertices has forest
proportion less than $\frac35$ for every even set $T$ of boundary vertices. The recursion is
sound, because $F(k)$ uses $F(k')$ with $k'<k$ only.

\paragraph{The base, $k<k_0$.} Here $2s_k\le2s_{k_0}<D_*$, so $N$ has no cycle of length at most
$2s_k$, and (f) with $t=0$ gives $P(k)$. If $N$ has at most one boundary vertex, then $T$ is
empty and $N$ contains a cycle, so the forest proportion is at most $\half$.

\paragraph{The step, $k\ge k_0$.} Let $\mathcal D$ be a maximal family of pairwise disjoint
cycles of length at most $2s$, and let $t$ be their number.
\begin{enumerate}[leftmargin=2.1em]
\item[(1)] If $t<k+n_1+n_2+2$, then $F(k)>3\cdot2^{s}(k+2s\,t)$, and (f) gives the claim.
\item[(2)] Let $t\ge k+n_1+n_2+2$. At most $k$ cycles of $\mathcal D$ contain a boundary vertex.
Each other cycle $Q$ has at most $2s$ boundary edges, and $V(Q)$ contains no vertex of $T$. So
more than $n_1$ of them are $T$-free, or more than $n_2$ of them are not.
\item[(3)] \emph{More than $n_1$ free cycles.} By (d) with $d=2s$, the forest proportion is less
than $\half+(2s+1)2^{2s-1}/n_1=\frac35$.
\item[(4)] \emph{More than $n_2$ cycles that are not free.} For such a cycle $Q$, some component
of $N-V(Q)$ contains an odd number of vertices of $T$. As $|T|$ is even and $T$ avoids $Q$, there
are $a,b\in T$ in different components of $N-V(Q)$. There are fewer than $k^{2}$ pairs, so more
than $(k+2)w$ of these cycles separate the same two vertices $a$ and $b$. Their vertex sets are
pairwise disjoint and connected, and they avoid $a$ and $b$. By (c), with $u=a$ and $p=b$, they
are linearly ordered, and there are $k+2$ pairwise disjoint sets $Y_0,\dots,Y_{k+1}$.
\item[(5)] \emph{Two sets.} At most $k$ of the sets $Y_l$ contain a boundary vertex of $N$. So
there are two indices $l_1\ne l_2$ such that $Y_1:=Y_{l_1}$ and $Y_2:=Y_{l_2}$ contain none. These
two sets are disjoint. Their vertices have degree three in $N$, and they are not in $T$.
\item[(6)] \emph{Size.} Each $Y_i$ contains $w-1$ of the cycles, and a cycle has at least $D_*$
vertices. So $|Y_i|\ge(w-1)D_*\ge k'F(k')$.
\item[(7)] \emph{Boundary.} An edge that leaves $Y_i$ ends in one of two of the cycles, and each
of them has at most $2s$ boundary edges. So at most $4s=k'$ edges leave $Y_i$. The set $Y_i$ does
not contain $b$, and $N$ is connected. So every component of $N[Y_i]$ has an edge that leaves
$Y_i$, and there are at most $k'$ components.
\item[(8)] \emph{The two regions.} Let $Y_i'$ be the vertex set of a largest component of
$N[Y_i]$. Then $N[Y_i']$ is connected, $|Y_i'|\ge F(k')$, and $Y_1'\cap Y_2'=\emptyset$. A
vertex of $Y_i'$ has degree less than three in $N[Y_i']$ only if it is an end of an edge that
leaves $Y_i$. So $N[Y_i']$ has at most $k'$ boundary vertices. Its girth is at least $D_*$.
\item[(9)] \emph{Induction.} By $P(k')$, every even set of boundary vertices of $N[Y_i']$ has
forest proportion less than $\frac35$, for $i=1,2$.
\item[(10)] \emph{Product.} The sets $U_1=Y_1'$ and $U_2=Y_2'$ satisfy the hypotheses of (e):
they are disjoint, they avoid $T$, their vertices have degree three in $N$, and they induce
connected multigraphs. So the forest proportion of $T$ is at most $\frac9{25}<\frac35$.
\end{enumerate}

\paragraph{The size of $F$.} We have $2^{s}=O(k^{2}\log k)$ and $s=O(\log k)$, so
$F(k)=O\bigl(k^{6}\log^{5}k+k^{5}\log^{3}k\cdot F(4s_k)\bigr)$ for $k\ge k_0$. Let $k_1\ge k_0$
be such that this bound, with $F(4s_k)\le C(4s_k)^{7}$, is at most $Ck^{7}$ for every $C\ge1$
and $k\ge k_1$. Then $C_1=\max_{k<k_1}F(k)$ serves. \qed

\section{Rigid sets at girth five}\label{sec:rigid}

The ladder with $n$ rungs has girth four, four boundary vertices and a rigid set: the two ends
of one rail. So a multigraph of girth four with a rigid set can be large. At girth five it
cannot.

\begin{lemma}[Full pairs at girth five]\label{lem:fullfive}
Let $M$ be a connected subcubic multigraph of girth at least five. Let $x\ne y$ be vertices of
degree at most two such that every $\{x,y\}$-join of $M$ is a path. Then $|V(M)|<15\,d(M)$.
\end{lemma}

\begin{proof}
Put $d=d(M)$ and $b=\beta(M)$. If $b=0$, then $M$ is a tree and $|V(M)|=d-2$. Let $b\ge1$.

\step{A pole.} Delete a vertex of degree one other than $x$ and $y$, as long as there is one.
Such a vertex is on no join, so the joins and the rank do not change. If then $x$ has degree
one, every join contains the edge at $x$. Delete $x$ and take its neighbour $x'$ in the place of
$x$. The joins for $x'$ and $y$ are the old joins without that edge, so they are paths, and the
rank is the same. The vertex $x'$ is not $y$: otherwise the old joins are the edge $xy$ together
with an even subgraph, all even subgraphs are empty, and $b=0$. Do the same for $y$, and
repeat. The result $M_1$ is a pole in which the two vertices are ports and form a full pair, and
$\beta(M_1)=b$.

\step{The ports.} Let $u$ be a port of $M_1$. If $u$ has degree three in $M$, one of its edges
leads to a deleted tree. The vertices of this tree have degrees in $M$ that add up to twice the
number of its edges plus one. So the tree has a deficit of at least two. Different ports have
different trees. Hence $M_1$ has at most $d$ ports.

\step{The ladders.} Let $M_1^{\circ}$ and $N_1,\dots,N_q$ be as in~(b). An edge of
$M_1^{\circ}$ is a path of $M_1$ whose inner vertices are ports. Call it \emph{plain} if this
path has no inner vertex. At most $d-2$ edges are not plain. A cycle of $M_1^{\circ}$ with at
most four edges, all plain, would be a cycle of $M$. So every cycle of $M_1^{\circ}$ with at
most four edges has an edge that is not plain. Let $e_i$ be the number of edges of $N_i$ that
are not plain. If $N_i$ is a pair of parallel edges, then $e_i\ge1=b_i$. Let $N_i$ be a capped
ladder with $m$ tiles. Each square is a cycle of length four. Two consecutive single rungs and
the two edges between them form a cycle of length four. These cycles are pairwise
edge-disjoint, and there are at least $(m-1)/2$ of them, because a run of $j$ consecutive single
rungs gives $\lfloor j/2\rfloor$ cycles and there are at most one more run than squares. A cap
and the first tile contain a cycle of length three or four. So $e_i\ge\max\{1,(m-1)/2\}\ge m/3$.
The rank of $N_i$ is at most $2m+1\le6e_i+1\le7e_i$. Hence
\[
  b=\sum_ib_i\le7\sum_ie_i\le7(d-2),
\]
and~\eqref{eq:deficit} gives $|V(M)|=2b-2+d<15d$.
\end{proof}

\begin{lemma}[Rigid sets]\label{lem:rigidpiece}
Let $M$ be a connected subcubic multigraph of girth at least five with deficit $d$. Let
$\lambda\colon V(M)\to\mathbb Z$ be a map with $|\lambda(u)-\lambda(v)|\le1$ for every edge $uv$,
such that fewer than $W$ vertices have any given value. If some even set $U$ of boundary vertices
is rigid, then $|V(M)|\le500\,dW(W+d)$.
\end{lemma}

\begin{proof}
If $U=\emptyset$, every even subgraph is a forest, so $M$ has no cycle and $|V(M)|=d-2$. Let
$U\ne\emptyset$.

\step{A tree.} Let $S$ be a connected subgraph of $M$ that contains $U$ and has as few edges as
possible. It is a tree, and its vertices of degree one are in $U$. Call a vertex of $S$
\emph{special} if it is in $U$ or has degree three in $S$. There are at most $2|U|-2$ special
vertices. A \emph{segment} is a path of $S$ between two special vertices without a special inner
vertex. The tree $S$ is the union of at most $2|U|-3$ segments.

\step{A region on a long segment.} Put $\rho=30W+8d$ and $\ell_0=(2\rho+1)W+2$. Suppose that
some segment $\sigma$ has at least $2\ell_0$ edges, and let $z_0$ be a vertex of $\sigma$ with at
least $\ell_0$ edges of $\sigma$ on each side. Let $B$ be the vertex set of the component of
$z_0$ in the subgraph of $M$ induced by the vertices $v$ with $|\lambda(v)-\lambda(z_0)|\le\rho$.
Then $M[B]$ is connected, and $B$ has fewer than $(2\rho+1)W=\ell_0-2$ vertices. The values of
$\lambda$ change by at most one along an edge. So $B$ contains the $2\rho+1$ vertices of
$\sigma$ nearest to $z_0$.

\step{The tree meets $B$ in one part of $\sigma$.} Let $v$ be a vertex of $S$ in $B$, and let
$\pi$ be a path from $v$ to $z_0$ in $M[B]$. It has fewer than $\ell_0-2$ edges. Suppose that
$v$ is not an inner vertex of $\sigma$. Delete from $S$ the part of $\sigma$ between $z_0$ and
the end of $\sigma$ on the side of $v$, and add $\pi$. The result is connected, it contains $U$,
and it has fewer edges, because the deleted part has at least $\ell_0$ edges. This is not
possible. In the same way, if $v$ is an inner vertex of $\sigma$, then $\sigma$ has at most as
many edges between $v$ and $z_0$ as $\pi$ has. So the vertices of $S$ in $B$ are inner vertices
of $\sigma$, and $B$ contains no vertex of $U$.

Let $z'$ and $z''$ be the ends of $\sigma$, and let $S^{-}$ and $S^{+}$ be the components of
$S-B$ that contain them. They are different, because the path of $S$ from $z'$ to $z''$ is
$\sigma$, which contains $z_0$. Every special vertex is joined to $z'$ or to $z''$ by a path of
$S$ without inner vertices of $\sigma$. So every vertex of $U$ is in $S^{-}$ or in $S^{+}$. Let
$x^{-}$ be the first vertex of $\sigma$ in $B$, seen from $z'$, and let $x^{+}$ be the first one
seen from $z''$. Then $x^{-}\ne x^{+}$, the vertex $x^{-}$ has a neighbour in $S^{-}$, and
$x^{+}$ has a neighbour in $S^{+}$. Let $K^{-}$ and $K^{+}$ be the components of $M-B$ that
contain $S^{-}$ and $S^{+}$, and put $U^{-}=U\cap V(S^{-})$.

\step{The set $B$ is small.} We claim that $|B|<15\,d(M[B])$.

\emph{First case: $K^{-}=K^{+}$, or $|U^{-}|$ is even.} Then $M-B$ contains a $U$-join: in the
connected set $K^{-}=K^{+}$, or in $K^{-}$ for $U^{-}$ and in $K^{+}$ for $U\sm U^{-}$. Its
union with a cycle of $M[B]$ would be a $U$-join that is not a forest. So $M[B]$ is a tree, and
$|B|=d(M[B])-2$.

\emph{Second case: $K^{-}\ne K^{+}$ and $|U^{-}|$ is odd.} Let $J$ be an $\{x^{-},x^{+}\}$-join
of $M[B]$. Let $v^{-}$ be a neighbour of $x^{-}$ in $K^{-}$ and $v^{+}$ a neighbour of $x^{+}$
in $K^{+}$. The set $U^{-}$ with $v^{-}$ added or removed is even, so $K^{-}$ has a join for it;
in the same way $K^{+}$ has a join for $U\sm U^{-}$ with $v^{+}$ added or removed. These two
joins, the edges $x^{-}v^{-}$ and $x^{+}v^{+}$, and $J$ form a $U$-join of $M$. It is a forest,
so $J$ is a forest, and hence a path. The vertices $x^{-}$ and $x^{+}$ have degree at most two
in $M[B]$. So the claim is Lemma~\ref{lem:fullfive} for $M[B]$.

An edge of $M$ that leaves $B$ starts in a vertex $v$ with $|\lambda(v)-\lambda(z_0)|=\rho$,
because $B$ is a component. There are fewer than $2W$ such vertices, and each has at most two
such edges. So $d(M[B])<4W+d$, and the claim gives $|B|<60W+15d<2\rho+1$. This contradicts
$|B|\ge2\rho+1$. So every segment has fewer than $2\ell_0$ edges.

\step{The count.} So $S$ has fewer than $4|U|\,\ell_0\le4d\,\ell_0$ vertices. The set $V(S)$ is
connected and contains $U$. So $M-V(S)$ is a forest, as in the first case. Let it have $f$
vertices. At most $|V(S)|+2$ edges join it to $S$, and a count of the edge ends in the forest
gives $f\le|V(S)|+d$. Hence
$|V(M)|\le8d\,\ell_0+d=8d\bigl((60W+16d+1)W+2\bigr)+d\le500\,dW(W+d)$.
\end{proof}

\begin{fineprint}
The proof has three ideas. A rigid set stays rigid in a part of the multigraph, for every way in
which the joins can cross the boundary of the part. A tree that connects the terminals with the
fewest edges passes through a region around a far point of the tree one time, so this region
has exactly two sides with terminals. If both sides are odd, one strand passes, and the region
has a full pair; if they are even, a join can avoid the region, and the region has no cycle.
\end{fineprint}

\section{Pieces at girth five}\label{sec:pieces}

\begin{theorem}[Pieces with few boundary vertices]\label{thm:base}
For every integer $k\ge1$ put $W=B(k+1)^{2}$, $h=10^{4}\,W^{3}$ and
$F_0(k)=Wh\,(k+2^{Wh}+2)$. Let $N$ be a connected subcubic multigraph of girth at least five with
at most $k$ boundary vertices and at least $F_0(k)$ vertices. Then every set of boundary vertices
of even size has forest proportion less than $\frac35$.
\end{theorem}

\begin{proof}
Suppose that some even set $T$ of boundary vertices has forest proportion at least $\frac35$.
If $T$ is empty, then $N$ has no cycle. A tree with at most $k$ boundary vertices has fewer than
$2k$ vertices. So let $T\ne\emptyset$. Fix a boundary vertex $o$, and call the set of the
vertices at distance $j$ from $o$ the layer $j$.

\step{Thin.} By (a), each layer has fewer than $W$ vertices. So $N$ has at least
$h\,(k+2^{Wh}+2)$ layers.

\step{Stretches.} Divide the layers $1,2,\dots$ into groups of $h$ consecutive layers. There are
at least $k+2^{Wh}$ groups all of whose layers are not empty, and at least $2^{Wh}$ of them
contain no boundary vertex. Let $\Gamma$ be the union of the layers of such a group, and let $v$
be a vertex in its last layer. Let $M$ be the component of $N[\Gamma]$ that contains $v$. A
shortest path from $o$ to $v$ has one vertex in each layer, and its vertices in $\Gamma$ are in
$M$. So $|V(M)|\ge h$. The vertices of $M$ have degree three in $N$. An edge of $N$ that leaves
$M$ ends outside $\Gamma$, so it starts in the first or in the last layer of the group; and each
vertex of $M$ has a neighbour in $M$. So $d(M)<4W$. By~\eqref{eq:deficit} the rank of $M$ is
less than $Wh$.

\step{Not rigid.} The distance from $o$ in $N$ is a map $\lambda$ on $V(M)$ as in
Lemma~\ref{lem:rigidpiece}. If some even set of boundary vertices of $M$ were rigid, the lemma
would give $|V(M)|<500\cdot4W\cdot W\cdot5W=h$. So for every even set of boundary vertices of
$M$, at least one of the fewer than $2^{Wh}$ joins is not a forest, and the forest proportion is
at most $1-2^{-Wh}$.

\step{Product.} The sets $V(M)$ for the different groups are pairwise disjoint, they avoid $T$,
and their vertices have degree three in $N$. By (e), the forest proportion of $T$ is at most
$(1-2^{-Wh})^{2^{Wh}}<\half$. This is a contradiction.
\end{proof}

\begin{theorem}[Pieces at girth five]\label{thm:poly}
There is an absolute constant $C_1$ with the following property. Let $k\ge1$, and let $N$ be a
connected subcubic multigraph of girth at least five with at most $k$ boundary vertices and at
least $C_1k^{7}$ vertices. Then every set of boundary vertices of even size has forest
proportion less than $\frac35$.
\end{theorem}

\begin{proof}
Let $k_0$ be as in Section~\ref{sec:audit}. For $k<k_0$ put $F(k)=F_0(k)$. For $k\ge k_0$ define
$F(k)$ as in Section~\ref{sec:audit}, with $5$ in the place of $D_*$ in the definition of $w$.
The claim $P(k)$ is the statement of the theorem with $F(k)$ in the place of $C_1k^{7}$. For
$k<k_0$ it is Theorem~\ref{thm:base}. For $k\ge k_0$ the step of Section~\ref{sec:audit} holds
word for word: (1) uses (f), which has no hypothesis on the girth; (6) holds because a cycle has
at least five vertices; and in (8) and (9) the girth of $N[Y_i']$ is at least five. The bound
for $F(k)$ is the same, with $C_1=\max_{k<k_1}F(k)$.
\end{proof}

\begin{corollary}[Two regions]\label{cor:two}
Let $G$ be a cubic multigraph, and let $P$ be a set of vertices such that $G-P$ has girth at least
five. Suppose that $V(G)\sm P$ contains two disjoint sets $Y_1,Y_2$ and that, for some
$k_1,k_2\ge1$, each $G[Y_i]$ is connected, $|\delta Y_i|\le k_i$ and $|Y_i|\ge C_1k_i^{7}$. Then
$q(G,P)\le\frac9{25}$.
\end{corollary}

\begin{proof}
Theorem~\ref{thm:poly} and (e) with $T=\emptyset$; the proof of (e) does not change when the
forest is in $G-P$.
\end{proof}

\paragraph{The size of the constant.} For $k<k_0$ the number $F_0(k)$ is about
$2^{10^{4}B^{4}(k+1)^{8}}$, and $B>2^{100}$. So $C_1$ is a very large absolute constant. The
proof of \M{Theorem~8.1} needs a bound $F(k)\le2^{\varepsilon k}$ for the pieces with
$k\approx\log r$ boundary vertices; Theorem~\ref{thm:base} alone does not give this, and the
step of Section~\ref{sec:audit} is necessary for it.

\section{The last case of [M] without bounded girth}\label{sec:last}

\begin{theorem}[{\M{Theorem~8.1} for girth five}]\label{thm:half}
There is an integer $r_5$ with the following property. Let $G$ be a connected cubic multigraph of
rank $r\ge r_5$. Let $P$ be a nonempty set of at most $3\log r$ vertices such that $G[P]$ is
connected. If $G-P$ has girth at least five, then $q(G,P)<\frac25$.
\end{theorem}

\begin{proof}
The proof is the second form of~[N1] (\N{Section~4.3}), with Theorem~\ref{thm:poly} in the
place of \N{Theorem~4.4} and with $5$ in the place of $D_*$. I checked the four applications.

\emph{Blocks and the chain.} Let $A$, $u$, $p$ and $C_1,\dots,C_T$ be as in \N{Proposition~2.1},
with $T>(\Lambda+2)\bigl(\lceil2^{8}C_1\Lambda^{8}/m\rceil+1\bigr)$. The sets $Y_l$ of (c) avoid
$P$. At most $\Lambda$ of them contain a vertex with a neighbour in $A$, and two others remain.
For such a set $Y$, at most $2\Lambda$ edges leave $Y$, and
$|Y|\ge2\Lambda\cdot C_1(2\Lambda)^{7}$. So a component of $G[Y]$ has at least
$C_1(2\Lambda)^{7}$ vertices and at most $2\Lambda$ boundary vertices, and it is a subgraph of
$G-P$, so its girth is at least five. Corollary~\ref{cor:two} gives $q(G,P)\le\frac9{25}$.

\emph{Heavy families.} The proof of \N{Proposition~2.2} uses the chain and \M{Lemma~8.6} only.
The bound is
$\mu\le\frac32K+3\Lambda\kappa\,(\Lambda+3)\bigl(\lceil2^{8}C_1\Lambda^{8}/m\rceil+1\bigr)$.

\emph{Packing.} Let $d\ge5$, and let $Q_1,\dots,Q_n$ be pairwise disjoint cycles of length at
most $d$ in $G-P$. Each has at least five vertices and at most $d$ boundary edges. With
$\Lambda=d$, $m=5$, $\kappa=1$ and $K=d+1$,
\[
  \mu\;\le\;\tfrac32(d+1)+3d\,(d+3)\bigl(\lceil2^{8}C_1d^{8}/5\rceil+1\bigr)\;\le\;2^{9}C_1\,d^{10},
  \qquad n\le2^{d+1}\mu\le2^{10}C_1\,2^{d}d^{10} .
\]

\emph{The final step.} With $D=\lceil x/16\rceil$ and $L=\lfloor Ax\rfloor$, the deleted set has
$|S|\le2^{10}C_1\,2^{D}D^{11}\le r/x^{2}$ vertices. \M{Proposition~3.5} has no hypothesis on the
girth of $G-P$. The family of the cycles of $\mathcal C$ gives
$\mu\le\frac32(1+\eta)+3L\eta\,(L+3)\bigl(\lceil2^{12}C_1L^{8}/x\rceil+1\bigr)=\frac32+o(1)$,
because $\eta=r^{-1/60}$ and $C_1$ is a constant. This contradicts $\mu>30$.
\end{proof}

\paragraph{The proof of the main theorem with four cases.} In \M{Section~9}, after the cases of
a loop, of a cyclic cut of size at most three, and of a cyclic cut of size four, the multigraph
$G$ has no loop and no cyclic cut of size at most four, and $r\ge9$. Let $Z$ be a cycle with at
most four edges. At most four edges leave $V(Z)$, and at most eight edges have an end in $V(Z)$.
So $G-V(Z)$ has rank at least $r-8>0$, and $\delta V(Z)$ is a cyclic cut of size at most four.
Hence $G$ has girth at least five. Choose eight vertices $P$ such that $G[P]$ is connected.
Then $G-P$ has girth at least five, $q(G,P)<\frac25$ by Theorem~\ref{thm:half}, and the count
of \M{Section~9} gives $c(G)<(\frac25+\frac1{32})2^{r}<2^{r-1}$. So
\[
  G\text{ has no cyclic cut of size at most four},\quad r\text{ large}
  \quad\Longrightarrow\quad c(G)<2^{r-1},
\]
and Case~4 of~[M], with a cycle shorter than $D_*$, is not a separate case.

\paragraph{What is not necessary any more.} I searched the source of~[M] for each label.
\begin{center}
\small
\begin{tabular}{@{}p{5.4cm}p{8.8cm}@{}}
\toprule
Statement of [M] & Where it is used\\
\midrule
Theorem 7.1 (bounded girth), all of Section~7; 52 typeset lines & Case~4 of Section~9 only.\\\addlinespace
Corollary 5.6 (an induced pole with few ports has not many paths) & The proof of Theorem~7.1
only.\\\addlinespace
Lemma 5.5 (four thin layers); with Corollary~5.6, 38 typeset lines & The proof of Corollary~5.6
only.\\\addlinespace
The number $D_*$, the fifth question of Figure~2, the sixth of the eight sentences & The
statement of the cases.\\
\bottomrule
\end{tabular}
\end{center}
These statements stay: Proposition~5.3 and Theorem~5.4 (Section~6 uses them), Corollary~3.7
(Section~6 uses it), and Lemmas~5.2 and~5.8, which Section~8 now uses as well.

\paragraph{What is new in the text.} Lemma~\ref{lem:fullfive}, Lemma~\ref{lem:rigidpiece} and
Theorem~\ref{thm:base}, with (e) for joins. There are two forms of Section~8.
\begin{itemize}[leftmargin=1.2em]
\item \emph{The clean form.} One lemma on pieces, Theorem~\ref{thm:poly}. The chain, the heavy
families and the packing bound are as in the second form of~[N1]; there is no packing number in
the statements and no induction in the packing bound.
\item \emph{The lean form.} \M{Lemma~8.2} stays, for pieces with at least $2\Lambda_0$ boundary
vertices and no short cycle. Theorem~\ref{thm:base} treats the pieces with fewer boundary
vertices. The rest is the first form of~[N1]: the chain counts the blocks with a short cycle,
and the packing bound uses induction on the length. A block needs
$\Phi(\Lambda)=C_1\Lambda^{4}$ vertices for $\Lambda\ge\Lambda_0$ and a constant number $C_2$ of
vertices for $\Lambda<\Lambda_0$. The lean form does not need (d), (f) and the step of
Section~\ref{sec:audit}.
\end{itemize}

\paragraph{The size, measured.} I applied each form to a copy of the paper and built it. The
count is in lines of the text of the PDF file.
\begin{center}
\small
\begin{tabular}{@{}lccccc@{}}
\toprule
 & Now & [N1], first & [N1], second & Girth five, lean & Girth five, clean\\
\midrule
Section 5 (poles) & 339 & 339 & 339 & 301 & 301\\
Section 7 (bounded girth) & 52 & 52 & 52 & -- & --\\
Section 8 (the last case) & 271 & 259 & 289 & 358 & 384\\
The paper, typeset lines & 1624 & 1615 & 1644 & 1617 & 1643\\
The paper, pages & 42 & 42 & 42 & 42 & 43\\
Cases of the main theorem & 5 & 5 & 5 & 4 & 4\\
Statements in Sections 5, 7, 8 & 20 & 20 & 21 & 20 & 21\\
\bottomrule
\end{tabular}
\end{center}
So the removal of the case of bounded girth costs no length in the lean form: the paper has 7
lines less than now, the same number of statements, one section less and one case less. The
last section is longer by 87 lines, and Sections~5 and~7 lose 90 lines. The mathematical gain is
in the structure: one argument treats every graph without a small cyclic cut. The total amount
of proof is about the same.

\section{Is the statement true? Tests}\label{sec:tests}

The task asks for a search for counterexamples. I found none. All numbers below are complete
enumerations of the joins of small graphs by the program \texttt{fpmax.c}.

\paragraph{The ladder.} At girth four the ladder is a counterexample of every size: the two ends
of one rail form a rigid set (computed up to 16 rungs). Its part in the proof is visible: a
ladder is the capped ladder of \M{Lemma~5.8}, and Lemma~\ref{lem:fullfive} fails for it because
its cycles of length four have plain edges only.

\paragraph{Families of girth at least five.} For each graph the table gives the largest forest
proportion over all nonempty even sets $T$ of boundary vertices.
\begin{center}
\small
\begin{tabular}{@{}p{6.1cm}cp{6.6cm}@{}}
\toprule
Family & $k$ & Largest forest proportion, by the number $n$ of vertices\\
\midrule
Segments of generalised Petersen graphs with step 2 (girth 5) & 6 &
$1$ ($n=12$), $0.78$ ($24$), $0.64$ ($32$), $0.59$ ($36$), $0.48$ ($44$)\\\addlinespace
The same with step 3 (girth 6) & 8 & $1$ ($16$), $0.94$ ($24$), $0.84$ ($32$), $0.73$ ($40$)\\\addlinespace
Rings of length 5 with two and three edges between them in turn (girth 5; cuts of size four) &
4, 5 & $1$ ($10$), $0.88$ ($20$), $0.81$ ($30$), $0.75$ ($40$)\\\addlinespace
Tubes: rings of length 6, 8, 10 with 3, 4, 5 edges between consecutive rings (girth 6) &
6, 8, 10 & length 6: $1$ ($18$), $0.84$ ($36$); length 8: $0.99$ ($24$), $0.76$ ($40$);
length 10: $1$ ($30$), $0.91$ ($40$)\\\addlinespace
Rings over the Heawood graph, 7 edges between them; three rings; sets $T$ of size at most 4 &
14 & $0.62$ ($42$)\\\addlinespace
Petersen, Heawood, Pappus, dodecahedron, Desargues, McGee, Tutte--Coxeter without one vertex &
3 & from $0.88$ ($9$) to $0.55$ ($19$) and $0.62$ ($29$)\\\addlinespace
Random cubic graphs of girth at least 5 without two or three vertices of a path & 4, 5 &
$0.97$ to $1$ ($n\le12$), $0.74$ to $0.89$ ($23$, $24$), $0.68$ to $0.81$ ($31$, $32$)\\
\bottomrule
\end{tabular}
\end{center}
In each family the proportion decreases when the graph grows. No graph in these runs with more
than 30 vertices has a rigid set. The largest rigid set that I saw at girth at least five is the
set of all ten boundary vertices in a tube of three rings of length ten, with 30 vertices and
deficit 10.

\paragraph{What the numbers say about the threshold.} A piece with 4 to 10 boundary vertices and
about 40 vertices can have a forest proportion above $\frac35$. So the true threshold is not
very small. The theorem gives the exponent $7$ with a very large constant. The tests are far
from the range of the theorem. They can show an error of the statement at small size, and they
show none. They cannot confirm the exponent.

\paragraph{A search for large rigid sets.} Rigid sets stay rigid when an edge or a vertex of
degree one is deleted. So every multigraph with a rigid set arises from one edge when we add
vertices of degree one and edges, through multigraphs with rigid sets. The script
\texttt{rigid\_search.py} takes random steps of this kind and keeps at most $k$ boundary vertices
at every step. At girth four it reaches the cap of the rank, $14$, for $k=4,5,6$; it finds
ladders. At girth five it finds rank $0$, $5$ and $8$ for $k=4,5,6$, with two seeds. This search
space is small, so the result is weak evidence.

\paragraph{The reduction.} The script \texttt{rigid\_reduction\_check.py} tests the two cases in
the proof of Lemma~\ref{lem:rigidpiece} for arbitrary connected sets $B$ without a terminal: on
290 rigid sets of small graphs and 7272 sets $B$, a set with only even sides induces a forest,
and a set with exactly two odd sides has full pairs between them. It also tests
Lemma~\ref{lem:fullfive} on the 11420 full pairs at girth five that occur; the largest ratio
$|V|/d$ is $1.4$.

\section{The questions of the task}\label{sec:answers}

\emph{1. The audit of Theorem 4.4.} Section~\ref{sec:audit}. The proof is valid, and the theorem
is a tool: Theorem~\ref{thm:poly} is its form for girth five.

\emph{2. The theorem for pieces of a graph without a small cyclic cut.} It is true, with the
exponent $a=7$ and an absolute constant (Theorem~\ref{thm:poly}). The piece $N$ need not inherit
any connectivity: the proof uses only that $N$ has girth at least five. I found no
counterexample. The two kinds of counterexample are different. A counterexample to the bound
$\frac35$ is a family of pieces with a bounded number of boundary vertices, of unbounded size,
with a forest proportion of at least $\frac35$; Theorem~\ref{thm:base} excludes it at girth
five, and the ladder is such a family at girth four. A counterexample to the polynomial
threshold is a family with a forest proportion of at least $\frac35$ whose size grows faster
than every power of $k$; the step of Section~\ref{sec:audit} excludes it. The pieces with about
40 vertices in Section~\ref{sec:tests} are neither: they are below the threshold.

\emph{3. Can the induction start at girth five?} Yes. The base is Theorem~\ref{thm:base}. A
large piece with a fixed number of boundary vertices and a high forest proportion is thin,
by~(a). A long thin piece contains a stretch with a rigid set, by the product. A rigid set gives
a region with a full pair (Lemma~\ref{lem:rigidpiece}), and this region is a capped ladder up to
its ports (Lemma~\ref{lem:fullfive}). So the only obstruction is the ladder, and it contradicts
girth five. No small cyclic cut is necessary for the contradiction. The proof does not use
\M{Corollary~5.6}. It uses the exact statement \M{Lemma~5.8} on full pairs, and the step from
all even boundary sets to one pair of ports is Lemma~\ref{lem:rigidpiece}.

\emph{4. Does this remove Section 7?} Yes (Section~\ref{sec:last}). Theorem~7.1, Lemma~5.5 and
Corollary~5.6 are not necessary, and Corollary~5.6 is used nowhere else.

\emph{5. If girth five is false.} It is not false. For girth four the ladder is the
obstruction, and a graph without a cyclic cut of size at most four has no cycle of length four.

\section{Limits}\label{sec:limits}

\begin{itemize}[leftmargin=1.2em]
\item \emph{The constant.} $C_1$ is the largest value of $F(k)$ over a range of small $k$, and
there $F(k)$ is about $2^{10^{4}W^{4}}$ with $W=B(k+1)^{2}$. This has no effect on the main
theorem, which holds for large rank only. The cause is the product in the proof of
Theorem~\ref{thm:base}: it takes $2^{Wh}$ stretches. I found no polynomial argument for few
boundary vertices.
\item \emph{The amount of proof.} The removal of Section~7 does not make the paper shorter by
much, because Lemma~\ref{lem:fullfive}, Lemma~\ref{lem:rigidpiece} and Theorem~\ref{thm:base}
replace it. The proof of Theorem~7.1 with its two dependencies has 90 typeset lines, and the new
statements have 87 lines in the lean form.
\item \emph{The formal proof.} The Lean proof has the case of bounded girth and pieces of large
girth. With a new Section~8 the text and the formal proof differ in the whole last case. The
main statement is the same.
\item \emph{Dependencies.} Section~8 then depends on \M{Lemmas~5.2 and~5.8}. Before, it depended
on Sections~2 to~4 only.
\item \emph{No second reader.} The proofs of Lemma~\ref{lem:fullfive},
Lemma~\ref{lem:rigidpiece} and Theorem~\ref{thm:base} are new. They are short enough to check
by hand, and I ask for this check before the manuscript changes.
\end{itemize}

\section{Files and checks}\label{sec:files}

All files are in \texttt{paper/minimal/}.
\begin{center}
\small
\begin{tabular}{@{}p{7.3cm}p{6.9cm}@{}}
\toprule
File & Content and result\\
\midrule
\texttt{research/apply\_girth\_five.py} & Applies the girth-five form to a copy of the paper:
\texttt{clean} or \texttt{lean}. It writes Section~8, removes Section~7, Lemma~5.5 and
Corollary~5.6, and changes Sections~1 and~9. Both copies build with no undefined reference and
no overfull line.\\\addlinespace
\texttt{research/trial-girth-five-lean.pdf}\newline\texttt{research/trial-girth-five-clean.pdf}
& The two copies. The last case starts on page~31.\\\addlinespace
\texttt{checks/girth5/fpmax.c} & Forest proportions of all even boundary sets of a small
graph, by enumeration of the joins.\\\addlinespace
\texttt{checks/girth5/families.py} & The table of Section~\ref{sec:tests}. Exit
code~0.\\\addlinespace
\texttt{checks/girth5/rigid\_search.py} & The search for rigid sets. Seeds 1 and 2, 6000 steps.\\\addlinespace
\texttt{checks/girth5/rigid\_reduction\_check.py} & The two cases of
Lemma~\ref{lem:rigidpiece} and the bound of Lemma~\ref{lem:fullfive} on small graphs. Exit
code~0.\\
\bottomrule
\end{tabular}
\end{center}

\paragraph{The choice.} The lean form is the smaller change and the shorter text, and it keeps
\M{Lemma~8.2}. The clean form has one lemma on pieces and simpler statements after it. Each of
them removes the case of bounded girth. My recommendation is the lean form, after a second
reader has checked Section~\ref{sec:rigid} and Theorem~\ref{thm:base}.

\begin{thebibliography}{9}\small
\bibitem[M]{M} KNT, Astra and Fable~5.2, \emph{The maximum number of cycles in a graph of large
rank}, draft of 12 November 2026, 42 pages, file
\texttt{paper/minimal/exact-maximum/exact-maximum.pdf}.
\bibitem[N1]{N1} Fable~5.2, \emph{The large-girth theorem with one family of short cycles},
research note of 12 November 2026, file \texttt{paper/minimal/research/one-scale.pdf}.
\end{thebibliography}

\end{document}
